
\chapter{\LaTeX 简介}
\intro{
	Every time I read a LaTeX document, I think, wow, this must be correct!
	
	\rightline{——Christos Papadimitriou}
}

\LaTeX{} 是一种基于 \TeX{} 的排版系统，由美国计算机学家莱斯利·兰伯特（Leslie Lamport）在 20 世纪 80 年代初期开发。

它并不是一个简单的字处理软件（如 Microsoft Word），而是一种标记语言。用户通过编写代码来描述文档的结构，由编译器（如 pdfLaTeX 或 XeLaTeX）将其转化为精美的 PDF 文档。

缺点很显然:这个东西像编程语言一样，不够直观。但是Word等软件的排版效果和公式编辑器的功能远远不如 \LaTeX{}，尤其是在处理复杂的数学公式和长文档时。
于是人们发明了Markdown等更为简洁的标记语言，但它们在专业排版方面仍然无法与 \LaTeX{} 相媲美。
好吧其实如果这个的标题不是Latex的话，我认为Markdown这种轻排版用起来会更方便些。好吧咱们继续往下说。
\section{为什么选择 \LaTeX}
\begin{itemize}
    \item \textbf{专业性}：其生成的数学公式和排版效果达到出版级水准。
    \item \textbf{效率性}：自动处理目录、参考文献、公式编号和交叉引用。
    \item \textbf{稳定性}：在处理长达数百页的论文或书籍时，不会出现卡顿或格式错乱。
\end{itemize}

\section{数学公式示例}
在 \LaTeX{} 中，插入行内公式使用 \texttt{\$ ... \$}，例如：$E=mc^2$。

如果是独立行公式，可以使用 \texttt{\$\$ ... \$\$} 或 \texttt{equation} 环境：
\begin{equation}
    \int_{a}^{b} f(x) \,dx = F(b) - F(a)
\end{equation}
\section{\LaTeX 的安装和配置}
\subsection{安装网址和步骤}
\subsubsection{安装网址}
目前主流可能会使用 TeX Live 或 MiKTeX 作为 \LaTeX{} 的发行版。但有的时候老师会给你发Ctex。（因为我国际大学生物理学术竞赛的时候老师就给我发的这玩意）

别用CTex！因为这玩意已经不更新了！而且它的安装包里包含了很多过时的软件和宏包，可能会导致兼容性问题。建议直接使用 TeX Live 或 MiKTeX，这两者都在不断更新，支持最新的 \LaTeX{} 功能和宏包。

我自己用的Texlive，这个安装包比较大，所用的时间比较长（就是上次下载了20个小时的那个玩意，但是我同学说清华源只要半个小时）

下面是清华源的地址：https://mirrors.tuna.tsinghua.edu.cn/CTAN/systems/texlive/Images/

直接下载那个texlive.iso就OK。

\subsubsection{安装步骤}
下载完那个.iso文件后，使用解压软件(我是穷鬼，所以我直接用的360压缩)将其解压到一个文件夹中。然后右键以管理员身份运行解压后的文件夹中的安装程序install-tl-windows.bat，按照提示进行安装。

但有的时候我们会长时间停留在load界面（比如我），然后我从CSDN上面找了个方法：

\begin{enumerate}
    \item \textbf{启动终端：} 按下 \texttt{Win + R} 键，输入 \texttt{cmd} 并回车进入命令提示符。
    
    \item \textbf{切换目录：} 进入 TeX Live 镜像文件解压后的文件夹（假设在 G 盘）：
    \begin{lstlisting}
d:
cd G:\Texlive2025
    \end{lstlisting}
    
    \item \textbf{启动安装程序：} 输入以下命令以纯文本模式（非 GUI）启动安装：
    \begin{lstlisting}
install-tl-windows.bat --no-gui
    \end{lstlisting}
    
    \item \textbf{选择操作：} 
    \begin{itemize}
        \item 输入 \texttt{I}：按默认路径直接开始安装。
        \item 输入 \texttt{D}：进入路径设置界面。
    \end{itemize}

    \item \textbf{修改路径：} 若进入了 \texttt{D} 界面，输入数字 \texttt{1}，你会看到：
    \begin{quote}
        \texttt{New value for TEXDIR [C:/texlive/2025]:}
    \end{quote}
    在冒号后输入你自定义的安装路径（例如 \texttt{D:/Software/texlive/2025}）。

    \item \textbf{执行安装：} 路径修改完成后，输入 \texttt{R} 返回初始界面，最后输入 \texttt{I} 开始安装。
\end{enumerate}
之后的话按照上面提示走就可以了。

\subsection{选择编译器：pdfLaTeX vs. XeLaTeX vs. LuaLaTeX}
安装好 \LaTeX{} 发行版后，我们面对的第一个选择就是：用哪个编译器？

\begin{itemize}
    \item \textbf{pdfLaTeX}（命令：\texttt{pdflatex}）：最经典的编译器，速度最快。但它对 Unicode（中文等非英文字符）的支持是"打补丁"式的，处理中文需要用 \texttt{CJK} 宏包做额外配置。适合纯英文文档。
    
    \item \textbf{XeLaTeX}（命令：\texttt{xelatex}）\textbf{—— 推荐中文用户使用！}原生支持 Unicode，可以直接调用 Windows/macOS/Linux 系统中的任何字体（配合 \texttt{fontspec} 宏包），加上 \texttt{ctex} 宏包就能像 Word 一样自由切换中文字体。本书所有文档均使用 XeLaTeX 编译。
    
    \item \textbf{LuaLaTeX}（命令：\texttt{lualatex}）：在 XeLaTeX 基础上内嵌了 Lua 脚本语言，排版能力最强，但编译速度较慢，目前尚在过渡阶段。
\end{itemize}

以本书为例，因为使用了中文和系统字体，必须使用 \texttt{xelatex} 编译。对应到 VSCode 配置中，就是把 \texttt{xelatex} 放在 recipes 列表的第一个位置。

\subsection{配置VSCode}
安装好 \LaTeX{} 之后，我们还需要一个编辑器来编写 \LaTeX{} 文档。（就像我们写C++或者Python需要一个IDE一样）Visual Studio Code（VSCode） 是一个非常流行的代码编辑器，支持多种编程语言和扩展，非常适合用来编写 \LaTeX{} 文档。

这是VSCode的官网：https://code.visualstudio.com/

直接下载那个Stable版本就OK。（郁琰学长好像已经帮你下载过了？）

安装好VSCode之后，我们会发现居然是英文界面，吓死我了！不过不用担心，VSCode支持多语言界面，我们可以轻松切换到中文界面。下载一个插件就行，在左边那个三个方块组成的图标（扩展）里搜索 \textbf{Chinese (Simplified) Language Pack for Visual Studio Code}，安装后重启VSCode就可以了。

我们需要安装一个 \LaTeX{} 插件来支持 \LaTeX{} 的编写和编译。推荐使用 \textbf{LaTeX Workshop} 插件，它提供了丰富的功能，如语法高亮、自动补全、错误检查和 PDF 预览。
本人亲测十分好用！

下一步就是配置 \LaTeX{} 的编译器路径了。打开 VSCode 的设置（File -> Preferences -> Settings），搜索 \texttt{latex-workshop.latex.tools}，找到 \texttt{latex-workshop.latex.tools} 配置项，点击编辑（Edit in settings.json），在里面添加以下内容：
\begin{lstlisting}
{
    //%***************************** LaTeX Workshop 核心配置 *********************************%//
    "latex-workshop.showContextMenu": true,
    "latex-workshop.intellisense.package.enabled": true,
    "latex-workshop.latex.build.forceRecipeUsage": false,

    //%***************************** 编译指令 (Recipes) *********************************%//
    "latex-workshop.latex.recipes": [
        {
            "name": "xelatex",
            "tools": ["xelatex"]
        },
        {
            "name": "pdflatex",
            "tools": ["pdflatex"]
        },
        {
            "name": "xelatex -> bibtex -> xelatex*2",
            "tools": ["xelatex", "bibtex", "xelatex", "xelatex"]
        },
        {
            "name": "pdflatex -> bibtex -> pdflatex*2",
            "tools": ["pdflatex", "bibtex", "pdflatex", "pdflatex"]
        }
    ],

    //%***************************** 编译工具具体参数 (Tools) *********************************%//
    "latex-workshop.latex.tools": [
        {
            "name": "xelatex",
            "command": "xelatex",
            "args": [
                "-shell-escape",
                "-synctex=1",
                "-interaction=nonstopmode",
                "-file-line-error",
                "%DOCFILE%"
            ]
        },
        {
            "name": "pdflatex",
            "command": "pdflatex",
            "args": [
                "-shell-escape",
                "-synctex=1",
                "-interaction=nonstopmode",
                "-file-line-error",
                "%DOCFILE%"
            ]
        },
        {
            "name": "bibtex",
            "command": "bibtex",
            "args": ["%DOCFILE%"]
        }
    ],

    //%***************************** 预览与正反搜索配置 (内置查看器方案) *********************************%//
    // 使用 VS Code 内置标签页查看 PDF，不再依赖外部程序
    "latex-workshop.view.pdf.viewer": "tab",

    // 设置双击 PDF 即可跳回对应的代码行 (反向搜索)
    "latex-workshop.view.pdf.internal.synctex.keybinding": "double-click",

    // 每次保存文件时自动编译
    "latex-workshop.latex.autoBuild.run": "onSave",

    //%***************************** 辅助配置 *********************************%//
    "latex-workshop.latex.clean.fileTypes": [
        "*.aux", "*.bbl", "*.blg", "*.idx", "*.ind", "*.lof", "*.lot", 
        "*.out", "*.toc", "*.acn", "*.acr", "*.alg", "*.glg", "*.glo", 
        "*.gls", "*.ist", "*.fls", "*.log", "*.fdb_latexmk"
    ],
    "security.workspace.trust.untrustedFiles": "open",
    "explorer.confirmDelete": false
}
\end{lstlisting}
其中 \texttt{command} 字段需要根据你安装的 \LaTeX{} 发行版进行调整。如果你使用的是 TeX Live，通常命令就是 \texttt{xelatex} 和 \texttt{pdflatex}。如果你使用的是 MiKTeX，可能需要指定完整路径，例如 \texttt{C:/Program Files/MiKTeX/miktex/bin/xelatex.exe}。
配置完成后，保存 settings.json 文件。现在你就可以在 VSCode 中编写 \LaTeX{} 文档了！按下 \texttt{Ctrl + Alt + B} 就可以编译当前文档并预览生成的 PDF 了。

或者不想那么麻烦配置的话可以考虑下载Texstudio这个集成开发环境，它自带了编译器和预览功能，安装后就可以直接使用了。不过我个人觉得VSCode的插件生态更丰富，使用起来也更灵活一些。
下载官网https://texstudio.sourceforge.net/